\section{Why Law: Reason, Authority, Freedom, and the Common Good}

\subsection{Law as rational ordering}

Aquinas's compact definition names law as an ordinance of reason for the common good, made by one who has care of the community, and promulgated (\emph{Summa theologiae} I--II, \sourceurl{https://www.newadvent.org/summa/2090.htm}{q.~90, a.~4}). Each element excludes a deformation. Mere desire is not law because power wants it. Private advantage is not the public end of law. Good advice does not become juridically binding without competent authority. A secret norm cannot ordinarily coordinate a community as promulgated law.

This definition is analogical across divine and human law. God does not acquire authority from a political community, and natural law is not printed in an official gazette. Yet wisdom, order to good, authority, and sufficient manifestation remain the conceptual core. Human positive law makes the general and shared order practically determinate.

\subsection{Derivation by conclusion and determination}

Human law can derive from natural law in at least two ways (Aquinas, I--II, \sourceurl{https://www.newadvent.org/summa/2095.htm}{q.~95, a.~2}). A \emph{conclusion} follows from a more general moral premise: from the duty not to harm comes a prohibition of intentional killing. A \emph{determination} selects one reasonable specification among several: a community fixes a filing period, form of record, side of the road, monetary threshold, or competent office. Conclusions carry more of the natural norm's necessity; determinations carry the authority of the enactment within a field left open.

Most real laws mix both. A safety code serves life and bodily integrity, uses empirical judgments about risk, chooses measurable thresholds, distributes inspection duties, and creates procedures. Disagreement about one threshold need not deny natural law. Conversely, a legislator's freedom to choose a threshold does not license choices destructive of the end that gives the regulation its point.

\subsection{Why human law cannot prohibit every vice}

Human law addresses persons of differing virtue and must preserve a possible common order. Aquinas argues that it should not attempt to suppress every vice, but principally grave harms from which the community must be protected, while leading persons gradually toward virtue (I--II, \sourceurl{https://www.newadvent.org/summa/2096.htm}{q.~96, aa.~2--3}). That principle is neither moral relativism nor an automatic policy answer. It distinguishes the full field of morality from the subset that coercive public law can prudently regulate for a given community.

The distinction also protects the Church's legal order from a spiritualized maximalism. Canon law is ordered to ecclesial communion and salvation, but it does not convert every sin into a canonical delict or every counsel into an enforceable precept. Penal typicity, competence, procedure, proof, imputability, prescription, and defense of rights remain moral achievements, not obstacles to holiness.

\subsection{Natural law requires prudence, not intuitionism}

The first practical principle---good is to be done and pursued, evil avoided---is very general. Human inclinations toward life, family, truth, sociality, and the good known by reason supply basic directions, but concrete judgment requires facts, virtues, experience, and attention to circumstances. Aquinas expressly distinguishes common principles from conclusions that become increasingly contingent (I--II, \sourceurl{https://www.newadvent.org/summa/2094.htm}{q.~94, aa.~2, 4--6}). Passion, bad custom, social corruption, and mistaken reasoning can obscure even serious truths.

Therefore ``natural law says'' is the beginning of an argument, not its conclusion. A responsible claim identifies the relevant good, principle, inferential steps, facts, alternative specifications, authority, and degree of certainty. Revelation can confirm and clarify moral truth, but appeal to revelation in a plural civil forum does not remove the need to articulate publicly intelligible reasons when the claim is also said to be naturally knowable.

\subsection{Promulgation, conscience, and coercion}

Conscience is the practical judgment that an act is to be done or avoided. It is not a private legislature that makes good and evil, nor is it a mere feeling displaced whenever an official speaks. A person must follow a certain conscience and must form it truthfully. A mistaken conscience may alter subjective culpability without changing the objective good, the external juridical effects of an act, or the competence of a tribunal.

Positive law supplies public notice and institutional judgment, but moral obligation is not exhausted by enforceability. Natural and revealed duties can bind where no human court has competence. Conversely, a valid human rule can bind members of a community even when the underlying choice was not the only morally possible one. The obligation arises through legitimate authority's determination for a common good.

\subsection{Unjust law: distinguish five questions}

The maxim that an unjust law is ``not law'' expresses a defect in law's focal meaning. It must not be used as an improvised jurisdictional eraser. Aquinas distinguishes enactments contrary to human good through defective end, author, or distributive form from enactments opposed to the divine good; he also considers scandal and social disturbance (I--II, q.~96, a.~4). A complete analysis asks:

\begin{enumerate}
\item Is the enactment morally unjust, and with what certainty?
\item Does it lack moral power to bind in conscience, or does a residual duty remain to avoid grave disorder?
\item Is it legally invalid under the positive order's own constitution, or legally operative until set aside?
\item Does compliance constitute formal or material cooperation with wrongdoing, and are there proportionate reasons or less harmful alternatives?
\item What remedy is competent: interpretation, exception, recourse, litigation, legislative change, conscientious objection, civil disobedience, or refusal?
\end{enumerate}

John Paul II applies the tradition forcefully to laws attacking innocent life while still distinguishing direct participation, conscientious objection, and incremental harm reduction (\sourceurl{https://www.vatican.va/content/john-paul-ii/en/encyclicals/documents/hf_jp-ii_enc_25031995_evangelium-vitae.html}{\emph{Evangelium vitae} 68--74}). That application does not license a reader to transfer conclusions mechanically to every disputed regulation.
